We have analysed the dual of the SOS Lyapunov program for switched systems
and shown how to leverage it to study the system stability.
We also generalized the whole approach to \emph{constrained switched systems},
a class of systems that has attracted increasing attention recently.

It turns out from our analysis that these two concepts are intrinsically related:
Our \theoref{fourth}, which leverages the dual of the classical JSR algorithm,
actually naturally applies to the constrained case;
Even more, \propref{Bsc} transforms an unconstrained system into a scalar constrained one for the purpose of computing a lower bound.
Finally, we show in \theoref{lowrank} that unconstrained systems with low rank matrices naturally lead to the definition of an auxiliary \emph{constrained} system.


We have introduced two techniques to generate lower bounds from the solution of the SOS dual program.
In practice, these techniques provide periodic trajectories of high asymptotic growth rate.
Since the SOS program can be solved efficiently,
does this give an efficient algorithm to generate lower bounds on the CJSR with \emph{guaranteed accuracy}?
This is not clear, because our algorithm provides firm guarantees only when the computed measures are atomic, which is not always the case.

More generally, the techniques developed in this work, based on
generating ``bad'' trajectories for a dynamical system via dual solutions%
, naturally extend to many other problems in systems theory. We are currently
exploring such possibilities.
